%!TEX root = ../../main.tex
% Section 5: Conclusion

\section{Conclusion}\label{sec:conclusion}

We proposed the Autoencoding Factor Model (AFM), a neural autoencoder with a portfolio bottleneck that constrains each latent factor to be the return of an explicit, tradable portfolio. Portfolio weights and conditional factor loadings are produced by neural networks conditioned on mutual fund characteristics---economic sector holdings and asset class allocations derived from regulatory filings---and a learned market-state embedding from past returns. The entire system is trained end-to-end via reconstruction loss, jointly learning portfolio construction, factor loadings, and asset intercepts.

Trained on approximately 1,000 US mutual funds spanning 10 asset classes, the AFM with 5 factors achieves 85.0\% average out-of-sample $R^2$, a 16.6 percentage point improvement over the Fama--French 4-factor model. The largest gains arise on non-equity assets---US aggregate bonds (84\% vs.\ 11\%), emerging market debt (44\% vs.\ 11\%), and US high yield (53\% vs.\ 35\%)---where traditional equity-based factors have near-zero explanatory power. Ablation studies demonstrate that simpler architectures (single-layer MLP) outperform deeper alternatives (residual MLP) and that multi-asset training is critical for cross-asset generalization, enabling transfer learning of shared factor structure across equities, bonds, and REITs.

\paragraph{Limitations.} Several limitations warrant discussion. First, our evaluation is limited to explanatory $R^2$; we do not assess portfolio-level performance metrics such as Sharpe ratios, turnover, or transaction costs, which are essential for practical deployment. Second, municipal debt ($R^2 \approx 30\%$) and emerging market debt ($R^2 \approx 44\%$) remain challenging, suggesting that the current fund-level characteristics are insufficient for these asset classes---richer inputs such as credit quality, duration, and tax status may be needed. Third, results are based on a single train/validation split; cross-validation would strengthen robustness claims but is computationally prohibitive at present. Finally, the model lacks formal theoretical guarantees on factor recovery or out-of-sample generalization.

\paragraph{Future work.} Several directions are promising. Incorporating richer characteristics---credit ratings, duration profiles, macroeconomic indicators---may improve performance on fixed-income assets. Extending the framework to individual securities (stocks and bonds) rather than mutual funds would test scalability to larger cross-sections. A variational autoencoder (VAE) extension \citep{Duan2022factorvae} could provide uncertainty quantification over factor returns and loadings. Portfolio-level evaluation, including live trading simulation with realistic transaction costs, would bridge the gap between explanatory power and economic value. Finally, theoretical analysis of the portfolio bottleneck as an implicit regularizer---characterizing when and why it improves generalization relative to free latent codes---remains an open and important question.
